\documentclass[a4paper,12pt]{article}
\usepackage[utf8]{inputenc}
\usepackage[T1]{fontenc}
\usepackage[spanish]{babel}
\usepackage{amssymb}
\usepackage{geometry}
\geometry{margin=2.2cm}

\title{Checklist de la sesion 4}
\author{Entorno}
\date{}

\begin{document}

\maketitle

\section*{Python en VS Code}

Marca cada punto cuando lo hayas comprobado.

\begin{itemize}
    \item[$\square$] He creado la carpeta \texttt{e4}.
    \item[$\square$] He creado la carpeta \texttt{juegos}.
    \item[$\square$] He instalado la extension Python de Microsoft en VS Code.
    \item[$\square$] Entiendo que la extension de VS Code no sustituye al interprete de Python.
    \item[$\square$] He creado \texttt{probar\_python.py}.
    \item[$\square$] He comprobado Python con el boton Run Python File.
    \item[$\square$] Si Python no aparece, he pedido a Codex pasos desde la interfaz de VS Code.
    \item[$\square$] He creado \texttt{juegos/tres\_en\_raya.py}.
    \item[$\square$] He ejecutado el juego con Run Python File.
    \item[$\square$] Se ha abierto la ventana del juego.
    \item[$\square$] He jugado una partida de tres en raya.
    \item[$\square$] He pedido a Codex una explicacion sencilla del programa.
    \item[$\square$] He hecho un cambio pequeno en el codigo.
    \item[$\square$] He vuelto a ejecutar el juego y he comprobado el cambio.
    \item[$\square$] He creado la carpeta \texttt{entrega}.
    \item[$\square$] La carpeta \texttt{entrega} contiene \texttt{probar\_python.py} y \texttt{juegos/tres\_en\_raya.py}.
\end{itemize}

\section*{Nombre}

\vspace{1cm}

\hrule

\end{document}
